% Generated draft: requires mathematical and source review.
\documentclass[11pt]{article}
\usepackage[margin=25mm]{geometry}
\usepackage{fontspec}
\setmainfont{Latin Modern Roman}
\usepackage{amsmath,amssymb,mathtools,unicode-math}
\setmathfont{Latin Modern Math}
\usepackage{enumitem,needspace,xurl,hyperref}
\setlength{\parindent}{0pt}
\setlength{\parskip}{5pt}
\setlength{\emergencystretch}{3em}
\newcommand{\N}{\mathbb{N}}
\newcommand{\Z}{\mathbb{Z}}
\newcommand{\Q}{\mathbb{Q}}
\newcommand{\R}{\mathbb{R}}
\newcommand{\C}{\mathbb{C}}
\newcommand{\F}{\mathbb{F}}
\newcommand{\A}{\mathbb{A}}
\newcommand{\PP}{\mathbb{P}}
\newcommand{\E}{\mathbb{E}}
\newcommand{\eps}{\varepsilon}
\newcommand{\dd}{\,\mathrm{d}}
\begin{document}
\section*{UnsolvedMath ID: 800125}
\textbf{Source problem code:} catalogue entry
\subsection{Definitions and mathematical statement}
The following formulation uses the conventional smooth, unoriented Haefliger groupoid; the supplied extract does not specify regularity or orientation conventions, so these conventions require confirmation. Let \(q\) be the codimension. Let \(\Gamma_q\) be the groupoid of germs of smooth local diffeomorphisms of \(\R^q\), and let \(B\Gamma_q\) be its classifying space. The universal normal bundle determines a map
\[
\nu\colon B\Gamma_q\longrightarrow B\mathrm{GL}(q,\R).
\]
Write \(B\bar\Gamma_q\) for the homotopy fibre of \(\nu\) over the point representing the trivial bundle. Thus a framed codimension-\(q\) Haefliger structure has a classifying map to \(B\bar\Gamma_q\). A classifying map is called homotopically trivial if it is homotopic to a constant map.

A smooth codimension-\(q\) foliation \(\mathcal F\) on a smooth manifold \(M\) is Riemannian if its transverse tangent bundle admits a metric invariant under local holonomy. It is a Lie foliation if it admits transverse submersions to a fixed \(q\)-dimensional Lie group \(G\), with transition maps locally given by left translations. A choice of basis of the Lie algebra of \(G\) then induces a transverse framing. Write
\[
c_{\mathcal F}\colon M\longrightarrow B\Gamma_q
\]
for the underlying Haefliger classifying map, and, in the Lie case with this framing, write
\[
\bar c_{\mathcal F}\colon M\longrightarrow B\bar\Gamma_q
\]
for the framed classifying map.

The two numbered questions in the supplied record are:
\begin{enumerate}
\item Is \(c_{\mathcal F}\) null-homotopic for every codimension-\(q\) Riemannian foliation \(\mathcal F\)?
\item Is \(\bar c_{\mathcal F}\) null-homotopic for every codimension-\(q\) Lie foliation \(\mathcal F\)?
\end{enumerate}
The source extract gives no restrictions on the underlying manifold, such as compactness, connectedness, or the presence of boundary, and gives no explicit range for \(q\). The intended scope of these universal quantifiers therefore remains unresolved. No transverse orientability or framing assumption is explicitly supplied for the Riemannian question. Under the convention above, null-homotopy in \(B\Gamma_q\) would in particular imply triviality of the normal bundle; the intended conventions and any omitted hypotheses must be checked.

The extract also begins with an assertion of equivalence to the following question: is every \emph{(framed)} Haefliger structure whose microfoliation has the zero section as a leaf homotopically trivial? Here a microfoliation means a germ, along the zero section of a rank-\(q\) vector bundle \(E\to M\), of a codimension-\(q\) foliation transverse to the fibres. Framing means a specified trivialization of \(E\). Under the conventions above, the unframed and framed classifying maps take values in \(B\Gamma_q\) and \(B\bar\Gamma_q\), respectively. The antecedent of ``This amounts to asking'' is absent, so the relationship of this opening question to the two numbered questions cannot be asserted from the extract. The immediately following sentence reports a positive answer by Tsuboi for \(q=1\); its precise scope is unresolved.
\subsection{Short English statement}
The supplied Question 66 asks whether every Riemannian foliation has a null-homotopic Haefliger classifying map, and whether every Lie foliation has a null-homotopic framed Haefliger classifying map. The extract also mentions a zero-section-leaf formulation and a positive codimension-one result attributed to Tsuboi, but their relationship to the numbered questions is unclear. The registry labels the problem open; that label is supplied curation, not an independently verified assessment.
\subsection{Sources}
\begin{itemize}
\item Supplied local UnsolvedMath registry record: permanent ID \(800125\), exact source problem code \texttt{catalogue entry}, upstream ID \(800125\), upstream revision \texttt{d371ee917ac9647fbb06cd6fc529deec6f3dda82}. The supplied registry status is \texttt{open}.
\item Nariman and Yazdi, \emph{Problem List on Foliations and Diffeomorphism Groups} (2025), Question 66. Supplied source URL: \url{https://arxiv.org/abs/2503.18274}. The record identifies the extraction method as \texttt{arxiv\_latex} and gives the access date 2026-10-16. These are supplied provenance details; no independent source visit or verification is claimed.
\item Attribution supplied in the title and background: E. Ghys, appendix to the reference identified by the bibliography key \texttt{molino1988riemannian}. The supplied record does not include the full bibliographic entry or an appendix locator.
\item The extracted statement attributes a positive answer for \(q=1\) to Tsuboi. No supporting citation, date, precise theorem, or definitive identification of the question answered is supplied.
\end{itemize}
\end{document}
